% Introduction for paper 05 only (platform_hybrid/paper/paper.tex).
% Forked 2026-10-19 from sections/intro.tex, which is also consumed by the
% archived U01 compendium roots; edit this file, not the shared one.
\section{Introduction}
\label{sec:intro}
Comparing RL post-training systems requires separating the algorithm name from
the implemented objective, sampling protocol, checkpoint, and evaluator.
Implementation sensitivity is already established in policy-gradient
research~\citep{engstrom2020implementation,henderson2018deep}; heterogeneous
observations alone do not identify which implementation choice caused a gap.

This manuscript is a bounded audit of one run corpus. Its contributions are:
\begin{enumerate}
\item An objective-scoped account of why homogeneous rewards eliminate the
centred reward-contrast term in the specified noncanonical runner
(Section~\ref{sec:diagnostic-method}).
\item An artifact-bound inventory of the two reviewed synthetic tool-use cells,
with an immutable manifest and the alternative explanations it cannot exclude
(Section~\ref{sec:evidence-identity}).
\item A re-audit of the training-reward ledger under one standard: every
retained number is bound to a named artifact and a correctly labelled window,
or it is quarantined or withdrawn (Section~\ref{sec:main_results}).
\end{enumerate}
These contributions do not include a new diagnostic definition, a controlled
cross-library benchmark, a validated predictor, triage rule, or stopping rule,
or an algorithm ranking.

\paragraph{Relation to existing zero-advantage diagnostics.}
The fraction of groups with zero reward variance is already an operational
quantity in group-relative RL: AVSPO's Advantage Collapse Rate measures the
same object up to a tolerance~\citep{he2026advantage}, and dynamic sampling in
DAPO~\citep{yu2025dapo} and variance- or pass-rate-based prompt selection
(e.g., VCRL~\citep{jiang2025vcrl}) act on it. We therefore use the name
Zero-Variance Fraction (ZVF) only as a reporting field and claim no new
measurable, estimator, or intervention. Its complement is called the
\emph{mixed-group fraction} here (historical alias: Gradient Utilization, GU);
it measures only the fraction of groups with nonconstant rewards, not gradient
magnitude or usefulness. For binary rewards it should be reported signed
(all-wrong, all-correct, mixed) rather than as one number.

\paragraph{Why an audit.}
The corpus was assembled for an earlier, broader manuscript. Re-checking its
displays against the stored artifacts showed mislabelled windows, rows whose
two artifact versions disagree, and summaries with no released source
(Section~\ref{sec:main_results}). Readers of any text that reuses this corpus
need those dispositions; the relationship to those texts is mapped in
Section~\ref{sec:overlap}.
